\documentclass{article}
\usepackage{graphicx} % Required for inserting images
\usepackage{amsmath}
  \RequirePackage[french]{babel}
  \RequirePackage[T1]{fontenc}
  \RequirePackage[utf8]{inputenc}

\title{STT-2902 - Modélisation statistique}
\author{Julien Miron}


\begin{document}

\section*{Équipe 2 – Estimation ponctuelle}

\textbf{Session :} Automne 2026

\textbf{Style imposé : cours magistral.} Votre présentation doit prendre la forme d'un cours magistral structuré, comme si vous étiez les enseignants du cours. Ce style est obligatoire : votre présentation doit être rigoureuse, suivre un plan annoncé et permettre à l'auditoire de prendre des notes.

\vspace{0.5cm}
\textbf{Objectif :} Expliquer le principe de l'estimation ponctuelle : qu'est-ce qu'un estimateur, quelles sont ses propriétés souhaitables (biais, variance) et comment choisir entre plusieurs estimateurs. Vous n'avez pas à couvrir les sections 8.3, 8.4 et 8.5.

\vspace{0.5cm}
\textbf{Déroulement imposé} :
\begin{itemize}
    \item Un plan de cours annoncé au début et suivi tout au long de la présentation;
    \item Un développement théorique clair avec définitions, exemples et au moins une démonstration au tableau;
    \item Des exercices courts intégrés à la présentation (comme dans un vrai cours) avec un temps de réflexion pour l'auditoire;
    \item Un résumé des points clés à la fin.
\end{itemize}

\vspace{0.5cm}
\textbf{Requis} :
\begin{enumerate}
    \item Définir clairement la distinction entre paramètre, estimateur et estimation;
    \item Expliquer les notions de biais, de variance et d'erreur quadratique moyenne d'un estimateur, avec au moins un exemple d'estimateur biaisé;
    \item Présenter un exemple concret comparant deux estimateurs pour un même paramètre (ex. : moyenne vs médiane échantillonnale);
    \item Faire au moins une démonstration mathématique au tableau (ex. : montrer que $\bar{X}$ est un estimateur sans biais de $\mu$).
\end{enumerate}

\end{document}
